\documentclass[onecolumn,aps,prb,amsmath,amssymb,superscriptaddress,longbibliography]{revtex4-2}

\usepackage{newtxtext}
\usepackage{newtxmath}
\usepackage{microtype}
\usepackage{bm}
\usepackage{mathtools}
\usepackage{xcolor}
\definecolor{PRLBlue}{RGB}{46,48,146}
\usepackage[
    colorlinks,
    citecolor=PRLBlue,
    linkcolor=PRLBlue,
    urlcolor=PRLBlue,
]{hyperref}
\usepackage{orcidlink}

\newcommand{\Tr}{\operatorname{Tr}}
\newcommand{\ee}{\mathrm{e}}
\newcommand{\ii}{\mathrm{i}}
\newcommand{\trans}{\mathsf{T}}
\newcommand{\mc}[1]{\mathcal{#1}}
\newcommand{\bb}[1]{\bm{#1}}
\newcommand{\ror}[2]{\href{https://ror.org/#1}{#2}}

\allowdisplaybreaks

\begin{document}

\title{A Representation-Specific Sign-Free Mechanism in Fermionic Gaussian Quantum Monte Carlo}

\author{Yin-Kai Yu\,\orcidlink{0009-0006-2589-2772}}
\affiliation{Beijing National Laboratory for Condensed Matter Physics\\ \& \ror{05cvf7v30}{Institute of Physics, Chinese Academy of Sciences}, Beijing 100190, China}
\affiliation{\ror{05qbk4x57}{University of Chinese Academy of Sciences}, Beijing 100049, China}

\author{Zhong-Chao Wei\,\orcidlink{0000-0002-1708-4306}}
\email{zcwei@thp.uni-koeln.de}
\affiliation{Qingdao Binhai University, Qingdao 266555, China}

\author{Lei Wang\,\orcidlink{0009-0004-1348-8583}}
\email{wanglei@iphy.ac.cn}
\affiliation{Beijing National Laboratory for Condensed Matter Physics\\ \& \ror{05cvf7v30}{Institute of Physics, Chinese Academy of Sciences}, Beijing 100190, China}
\affiliation{\ror{05qbk4x57}{University of Chinese Academy of Sciences}, Beijing 100049, China}
\affiliation{\ror{020vtf184}{Songshan Lake Materials Laboratory}, Dongguan, Guangdong 523808, China}

\author{Zi-Xiang Li\,\orcidlink{0000-0002-3941-365X}}
\email{zixiangli@iphy.ac.cn}
\affiliation{Beijing National Laboratory for Condensed Matter Physics\\ \& \ror{05cvf7v30}{Institute of Physics, Chinese Academy of Sciences}, Beijing 100190, China}
\affiliation{\ror{05qbk4x57}{University of Chinese Academy of Sciences}, Beijing 100049, China}

\date{\today}

\begin{abstract}
Auxiliary-field quantum Monte Carlo weights each field configuration by an
ordered Fock-space trace of Gaussian propagators.  Majorana and Kramers time
reversals guarantee nonnegative traces from the symmetry algebra.  In six
fermionic relation classes, we find physical realizations whose complete
maximal-group spaces lift to commuting Hermitian Fock operators.  Every ordered
product reduces to one Hermitian exponential with a strictly positive trace for
arbitrary real fields.  In the two-mode realizations, the generators act in
complementary parity doublets and give
\(p_{\rm F}=2\cosh\alpha+2\cosh\beta\).  Every physical time reversal
preserving these spaces has positive square on the Nambu operators.  Positivity
therefore follows from parity-resolved commuting dynamics beyond Majorana and
Kramers pairing.  The analogous bosonic group direction has a pole between the
trace-convergent and negative continued branches, exposing the role of
Fock-space convergence and branch selection.
\end{abstract}

\maketitle

\tableofcontents

\section{Introduction}

Auxiliary-field quantum Monte Carlo transforms an interacting partition
function into a statistical sum over quadratic evolutions
selected by Hubbard--Stratonovich fields~\cite{blankenbecler1981prd}.  A
nonnegative Gaussian weight is the central requirement for direct importance
sampling.  Time-reversal symmetry supplies two general fermionic mechanisms.
The Majorana class separates the evolution into complex-conjugate sectors,
while the Kramers class enforces an even degeneracy of the relevant one-body
spectrum~\cite{li2015prb,li2016prl}.  Both mechanisms determine the sign from
the time-reversal algebra without reference to a particular matrix realization.

The companion note \textit{Enumeration of Contraction Semigroups for
Sign-Problem-Free Quantum Monte Carlo} organizes Gaussian evolutions through
two real orthogonal transformations \(J_1\) and \(J_2\).  Five binary
relations specify their abstract algebra~\cite{wei2024prb}.  The companion note
\textit{Time-Reversal Symmetries for Sign-Problem-Free Quantum Monte Carlo}
classifies the Majorana, Kramers, and bosonic split-orthogonal mechanisms.
These results are sufficient criteria.  A concrete physical realization also
fixes the dimension, the action of \(J_1,J_2\) on the Nambu operators, and the
complete algebra of allowed quadratic operators.  This additional information
can fix the sign even when the five relations alone do not.

Here we establish a representation-specific fermionic mechanism.  The term
representation-specific refers to physically different realizations of the
same abstract relations, including different numbers of fermion modes and
different actions of \(J_1,J_2\).  A unitary change of basis within one
realization leaves the Fock trace unchanged.  In the examples below, the
complete maximal-group generators become commuting Hermitian operators on a
finite fermion Fock space.  Their support separates into the even- and
odd-parity sectors.  This gives exact positive weights for arbitrary numbers
and orderings of time slices.

\section{Gaussian weights and a representation-specific criterion}

For one auxiliary-field configuration, bosons or fermions evolve through
quadratic operators
\begin{equation}
    p=\Tr\!\left(\ee^{\hat h_1}\ee^{\hat h_2}\cdots\ee^{\hat h_L}\right),
    \qquad
    \hat h_\ell=\frac12\hat\chi^\trans gA_\ell\hat\chi .
    \label{eq:configuration_weight}
\end{equation}
The Nambu column is
\begin{equation}
    \hat\chi=(\hat a_1,\ldots,\hat a_N,
    \hat a_1^+,\ldots,\hat a_N^+)^\trans ,
    \label{eq:nambu_column}
\end{equation}
with canonical matrices
\begin{equation}
    g_{\rm B}=\begin{pmatrix}0&-I_N\\ I_N&0\end{pmatrix},
    \qquad
    g_{\rm F}=\begin{pmatrix}0&I_N\\ I_N&0\end{pmatrix}.
    \label{eq:canonical_matrices}
\end{equation}
The quadratic condition
\begin{equation}
    gA_\ell+A_\ell^\trans g=0
    \label{eq:quadratic_condition}
\end{equation}
makes \(gA_\ell\) symmetric for bosons and antisymmetric for fermions.  The
ordered one-body propagation is
\begin{equation}
    M=\ee^{A_1}\ee^{A_2}\cdots\ee^{A_L},
    \label{eq:one_body_propagation}
\end{equation}
and the Gaussian trace identities are~\cite{munn1978ensemble,klich2014jstat}
\begin{equation}
    p_{\rm B}=\left[(-1)^N\det(I_{2N}-M)\right]^{-1/2},
    \qquad
    p_{\rm F}=\left[\det(I_{2N}+M)\right]^{1/2}.
    \label{eq:gaussian_trace_identities}
\end{equation}
The ordered Fock-space product fixes the overall square-root sign that is
absent from the determinant.

The contraction-semigroup classification imposes
\begin{equation}
\begin{gathered}
    J_1^\trans gJ_1=\pm g,\qquad J_1^2=\pm I,\\
    J_2^\trans gJ_2=\pm g,\qquad J_2^2=\pm I,\qquad
    J_1J_2=\pm J_2J_1 .
\end{gathered}
    \label{eq:five_relations}
\end{equation}
For a chosen realization of \(J_1,J_2\), the complete maximal-group generator
space is the real solution of
\begin{equation}
\begin{gathered}
    gA+A^\trans g=0,\\
    J_1^\trans AJ_1=\overline A,\qquad
    J_2^\trans AJ_2=\overline A .
\end{gathered}
    \label{eq:complete_group_space}
\end{equation}
Both \(A\) and \(-A\) belong to this real space.  The corresponding
propagations therefore include inverse elements, and auxiliary fields can take
either sign along a group direction.

The five signs in Eq.~\eqref{eq:five_relations} do not determine the dimension
or multiplication rules of this solution space.

Let \(A_a\) be a complete real basis of
Eq.~\eqref{eq:complete_group_space}, and let
\(\hat H_a=\frac12\hat\chi^\trans gA_a\hat\chi\) be its Fock-space lift.
Suppose
\begin{equation}
    \hat H_a^\dagger=\hat H_a,
    \qquad
    [\hat H_a,\hat H_b]=0
    \quad\text{for all }a,b .
    \label{eq:commuting_hermitian_condition}
\end{equation}
Every slice has the form
\(\hat h_\ell=\sum_a x_{\ell a}\hat H_a\) with real coefficients.  All slices
then commute, and Eq.~\eqref{eq:configuration_weight} becomes
\begin{equation}
\begin{aligned}
    p_{\rm F}
    &=\Tr\exp\!\left(\sum_{\ell,a}x_{\ell a}\hat H_a\right)\\
    &=\sum_n \exp(\lambda_n)>0 .
\end{aligned}
    \label{eq:fermion_positive_trace}
\end{equation}
The \(\lambda_n\) are real eigenvalues of the summed Fock-space generator.
The finite fermion Hilbert space makes every term finite.  This argument fixes
the sign for arbitrary real fields and for any number of time slices.

\section{Parity-resolved fermion realizations}

\subsection{One fermion mode}

For one fermion mode,
\(\hat\chi=(\hat c,\hat c^+)^\trans\) and \(g_{\rm F}=\sigma_x\).
Four choices of \(J_1,J_2\) realize four of the relation classes in the
contraction-semigroup note:
\begin{equation}
\begin{array}{c|c}
\text{class} & (J_1,J_2) \\ \hline
\mathrm{F01} & (I_2,I_2)\\
\mathrm{F05} & (I_2,\sigma_z)\\
\mathrm{F17} & (\sigma_z,I_2)\\
\mathrm{F21} & (\sigma_z,\sigma_z).
\end{array}
    \label{eq:one_mode_realizations}
\end{equation}
Direct substitution into Eq.~\eqref{eq:complete_group_space} gives the same
complete real space in all four cases,
\begin{equation}
    A_{\rm group}=a\,\sigma_z,\qquad a\in\mathbb R .
    \label{eq:one_mode_group}
\end{equation}
Its Fock-space lift is
\begin{equation}
    \hat h=a\left(\hat n-\frac12\right).
    \label{eq:one_mode_fermion_generator}
\end{equation}
For a sequence of slices, let \(\alpha=\sum_\ell a_\ell\).  The two occupation
states give
\begin{equation}
    p_{\rm F}
    =\Tr\ee^{\alpha(\hat n-1/2)}
    =\ee^{-\alpha/2}+\ee^{\alpha/2}
    =2\cosh\frac{\alpha}{2}>0 .
    \label{eq:one_mode_fermion_weight}
\end{equation}

\subsection{Two fermion modes}

The two-mode mechanism has a simple physical form.  In the occupation basis
\(\{|00\rangle,|01\rangle,|10\rangle,|11\rangle\}\), define the odd-parity
generator
\begin{equation}
    \hat H_{\rm o}
    =\hat c_1^+\hat c_2+\hat c_2^+\hat c_1 .
    \label{eq:odd_parity_generator}
\end{equation}
It exchanges \(|01\rangle\) and \(|10\rangle\) and annihilates both
even-parity states.  Depending on the realization, the second generator is one
of the following three Hermitian operators.  Each acts only in the
complementary even-parity doublet:
\begin{align}
    \hat H_{\rm e}^{(1)}
      &=\hat n_1+\hat n_2-1,\notag\\
    \hat H_{\rm e}^{(2)}
      &=\ii\left(\hat c_1\hat c_2+
        \hat c_1^+\hat c_2^+\right),\notag\\
    \hat H_{\rm e}^{(3)}
      &=\hat c_1^+\hat c_2^+-\hat c_1\hat c_2 .
    \label{eq:even_parity_generators}
\end{align}
The first is diagonal in \(\{|00\rangle,|11\rangle\}\).  The second and third
rotate this pair with different phase conventions.  Each has eigenvalues
\((-1,+1)\) in its active doublet.  The odd generator has the same two
eigenvalues in the odd doublet.  Each complete two-mode generator space is
spanned by \(\hat H_{\rm o}\) and one \(\hat H_{\rm e}^{(j)}\).  Their disjoint
support gives
\begin{equation}
    \hat H_{\rm e}^{(j)}\hat H_{\rm o}
    =\hat H_{\rm o}\hat H_{\rm e}^{(j)}=0 .
    \label{eq:parity_support}
\end{equation}

The F01 realization uses
\(\hat H_{\rm e}^{(1)},\hat H_{\rm o}\), the F06 and F18 realizations use
\(\hat H_{\rm e}^{(2)},\hat H_{\rm o}\), and the F21 realization uses
\(\hat H_{\rm e}^{(3)},\hat H_{\rm o}\).  Appendix~\ref{app:explicit_matrices}
gives the corresponding Nambu matrices and verifies that these pairs span the
complete spaces in Eq.~\eqref{eq:complete_group_space}.  For
\(\alpha=\sum_\ell a_\ell\) and
\(\beta=\sum_\ell b_\ell\), Eq.~\eqref{eq:parity_support} gives the exact
weight
\begin{equation}
\begin{aligned}
    p_{\rm F}
    &=\Tr\ee^{\alpha\hat H_{\rm e}^{(j)}
                  +\beta\hat H_{\rm o}}\\
    &=\underbrace{\ee^{-\alpha}+\ee^\alpha}_{\text{even parity}}
      +\underbrace{\ee^{-\beta}+\ee^\beta}_{\text{odd parity}}\\
    &=2\cosh\alpha+2\cosh\beta\geq4 .
\end{aligned}
    \label{eq:two_mode_fermion_weight}
\end{equation}
The one-mode and two-mode realizations together establish the
representation-specific mechanism in the six classes F01, F05, F06, F17,
F18, and F21.

\section{Relation to time-reversal positivity}

The Majorana and Kramers criteria require pairs of anticommuting physical time
reversals with square assignments \(\{+,-\}\) and \(\{-,-\}\),
respectively~\cite{li2016prl}.  Each pair contains a time reversal that squares
to \(-I\).  The complete generator spaces above admit no such member.

Let \(\widehat\Theta\) be an antiunitary transformation on Fock space, and let
\(\Theta(\hat d)=\widehat\Theta\hat d\widehat\Theta^{-1}\) denote its induced
anti-linear action on an operator.  On the Nambu vector,
\begin{equation}
    \Theta(\hat\chi)=U\hat\chi,
    \qquad
    \Theta(z\hat d)=\overline z\,\Theta(\hat d),
    \label{eq:time_reversal_action}
\end{equation}
or equivalently \(\Theta=U\mc K\), where \(\mc K\) conjugates numerical
coefficients.  The Nambu exchange matrix
\begin{equation}
    C=\begin{pmatrix}0&I_N\\ I_N&0\end{pmatrix}
    \label{eq:nambu_exchange}
\end{equation}
encodes Hermitian conjugation and equals \(g_{\rm F}\) for fermions.  A
physical time reversal satisfies
\begin{equation}
    U^\trans C U=C,\qquad
    C\overline U C=U,\qquad
    \overline A\,U=UA
    \quad\text{for every allowed }A .
    \label{eq:physical_antiunitary}
\end{equation}
Solving the adjoint and invariance equations for the complete one- and
two-mode spaces gives the matrix families listed in
Appendix~\ref{app:antiunitary_spaces}.  Every matrix in those families is
symmetric.

The physical constraints then fix the square without further case-by-case
analysis.  Transposing the second equation in
Eq.~\eqref{eq:physical_antiunitary} gives
\(U^\trans=CU^\dagger C\).  Substitution into the first equation yields
\begin{equation}
    U^\dagger U=I .
    \label{eq:antiunitary_unitarity}
\end{equation}
Symmetry gives \(U^\dagger=\overline U\).  Applying Eq.~\eqref{eq:time_reversal_action}
twice gives
\begin{equation}
    \Theta^2(\hat\chi)
    =\overline U U\hat\chi
    =U^\dagger U\hat\chi
    =\hat\chi .
    \label{eq:antiunitary_square}
\end{equation}
The complete invariant time-reversal family contains only positive-square
members.  Positive-square symmetries can accompany the evolution, while the
fermionic Majorana and Kramers mechanisms require a negative-square member.
The strict positivity in Eqs.~\eqref{eq:one_mode_fermion_weight} and
\eqref{eq:two_mode_fermion_weight} comes from the finite Fock spectrum and the
separation of fermion-parity sectors.

The contraction-semigroup note marks all six relation classes by a proper
time-reversal-protected subgroup.  This label records the subgroup reached by
the Majorana or Kramers sufficient criterion.  For F01, F06, F18, and F21, the
two-mode realizations above are the same realizations displayed there, and
their entire maximal groups are sign free.  For F05 and F17, the one-mode
realizations establish sign-free maximal groups in alternative physical
realizations of the same five relations.

\section{Bosonic convergence and branch selection}

The finite fermion Hilbert space is essential to
Eq.~\eqref{eq:fermion_positive_trace}.  A bosonic Fock space is infinite.  A
real group direction contains both a Hermitian generator \(\hat H\) and
\(-\hat H\).  Their exponentials cannot both have finite traces on an
infinite-dimensional Hilbert space.  Physically, the direction that damps
large occupation numbers becomes amplifying when its sign is reversed.
Equivalently, if both exponentials were trace class, their product would make
the identity trace class.

The simplest one-mode realization makes the branch structure explicit.  Let
\(\hat\chi=(\hat b,\hat b^+)^\trans\), \(g_{\rm B}=-\ii\sigma_y\), and choose
\begin{equation}
    (J_1,J_2)=(I_2,\sigma_z)
    \quad\text{or}\quad
    (\sigma_z,I_2).
    \label{eq:boson_J_pairs}
\end{equation}
These are the B05 and B17 relations in the contraction-semigroup note.
Solving Eq.~\eqref{eq:complete_group_space} gives
\begin{equation}
    A_{\rm group}=a\,\sigma_z,\qquad
    \hat h=a\left(\hat n+\frac12\right).
    \label{eq:boson_group_generator}
\end{equation}
For \(\alpha=\sum_\ell a_\ell<0\), the ordinary Fock trace converges and equals
\begin{equation}
\begin{aligned}
    p_{\rm B}
    &=\sum_{n=0}^\infty \ee^{\alpha(n+1/2)}
      =\frac{\ee^{\alpha/2}}{1-\ee^\alpha}\\
    &=\frac{1}{2\sinh(-\alpha/2)}>0 .
\end{aligned}
    \label{eq:boson_convergent_trace}
\end{equation}
The one-body propagation is \(M=\ee^{\alpha\sigma_z}\), so
\begin{equation}
    -\det(I_2-M)=4\sinh^2\frac{\alpha}{2}.
    \label{eq:boson_determinant}
\end{equation}
Within the pure one-parameter group, the meromorphic continuation of
Eq.~\eqref{eq:boson_convergent_trace} to \(\alpha>0\) is
\begin{equation}
    p_{\rm B}
    =-\frac{1}{2\sinh(\alpha/2)}<0 .
    \label{eq:boson_negative_branch}
\end{equation}
The occupation-number sum diverges there.  The positive determinant in
Eq.~\eqref{eq:boson_determinant} cannot distinguish the two square-root signs.

The contraction semigroup is specified by a Hermitian metric \(\eta\) with
\(\eta^2=I\).  Its strict interior obeys
\(\eta A+A^\dagger\eta<0\), which keeps the ordered propagation in a connected
pole-free region anchored at a positive reference~\cite{wei2024prb}.  For the
two choices in Eq.~\eqref{eq:boson_J_pairs}, take \(R>0\) and
\begin{equation}
\begin{array}{c|c|c}
(J_1,J_2) & \eta & A(\theta) \\ \hline
(I_2,\sigma_z)
&-\sigma_x
&R(\cos\theta\,\sigma_z+\sin\theta\,\sigma_x)\\
(\sigma_z,I_2)
&\sigma_y
&R(\cos\theta\,\sigma_z-\sin\theta\,\sigma_y).
\end{array}
    \label{eq:boson_contraction_paths}
\end{equation}
In each row, the term proportional to \(\sin\theta\) is a positive multiple of
the cone direction \(-\eta\).
For \(0<\theta<\pi\), both rows obey
\begin{equation}
    \eta A(\theta)+A(\theta)^\dagger\eta
    =-2R\sin\theta\,I_2<0,
    \qquad
    A(\theta)^2=R^2I_2 .
    \label{eq:boson_path_properties}
\end{equation}
The determinant radicand remains
\(4\sinh^2(R/2)\) along the entire path.  At \(\theta=\pi\), the generator is
\(-R\sigma_z\) and its ordinary Fock trace is finite and positive.  The
contraction inequality keeps the continued Gaussian weight finite and on this
positive branch as \(\theta\) moves to zero.  The endpoint is
\(+R\sigma_z\), where the ordinary trace diverges.

The same one-body endpoint therefore carries opposite signs under the pure
group continuation in Eq.~\eqref{eq:boson_negative_branch} and the
cone-anchored continuation in Eq.~\eqref{eq:boson_contraction_paths}.  The
ordered continuation path supplies information that the determinant has
lost.  The positive cone-anchored branch is precisely the contraction-semigroup
mechanism developed in the companion note.  Commuting Hermitian generators
alone consequently do not provide a bosonic full-group counterpart of the
fermionic mechanism.

\section{Summary}

The algebraic relations among anti-linear transformations give powerful
representation-independent sufficient conditions for sign-free Gaussian
quantum Monte Carlo.  A physical realization contains additional information.
It determines the full allowed generator space and how that space acts on the
many-body Hilbert space.

We identified realizations in six fermionic relation classes whose complete
maximal-group generators lift to commuting Hermitian Fock operators.  One-mode
realizations give \(2\cosh(\alpha/2)\).  Two-mode realizations divide the
dynamics between complementary even- and odd-parity doublets and give
\(2\cosh\alpha+2\cosh\beta\).  Every physical time reversal preserving these
spaces has positive square on the Nambu operators, so the positivity follows from
parity-resolved commuting dynamics beyond the Majorana and Kramers criteria.

Bosons expose the boundary of this mechanism.  Reversing a Hermitian group
direction reverses occupation-number damping, and an ordinary Fock trace
cannot converge on the complete real group.  The one-mode example further
shows that the pure-group and contraction-semigroup paths can select opposite
square-root signs at a nonconvergent endpoint.  Representation-specific
fermionic positivity and bosonic contraction-semigroup positivity therefore
arise from distinct physical structures.

\appendix

\section{Explicit Nambu realizations}
\label{app:explicit_matrices}

For the two-mode examples, write
\begin{equation}
    \hat\chi=(\hat c_1,\hat c_2,\hat c_1^+,\hat c_2^+)^\trans,
    \qquad
    g_{\rm F}=\sigma_x^{({\rm N})}\otimes I_2^{({\rm f})}.
    \label{eq:two_mode_nambu}
\end{equation}
The superscripts \({\rm N}\) and \({\rm f}\) label the Nambu and two-mode
factors.  The matrices \(\ii\sigma_y\) are real.  Overall signs of \(J_1,J_2\)
and of individual basis matrices do not change the constraints or their real
spans.

\begin{center}
\refstepcounter{table}\label{tab:explicit_realizations}
\begin{minipage}{0.98\linewidth}
\small
\textbf{TABLE \thetable.}\enspace
Two-mode realizations and their complete maximal-group generator spaces.  The
two basis matrices lift in the order
\((\hat H_{\rm e},\hat H_{\rm o})\) used in the main text.
\par\vspace{0.6\baselineskip}
\centering
\begin{ruledtabular}
\begin{tabular}{ccl}
Class & \((J_1,J_2)\) & Complete \(A_{\rm group}\) space \\ \hline
F01 &
\(\bigl(I_2^{({\rm N})}\otimes I_2^{({\rm f})},
I_2^{({\rm N})}\otimes\sigma_x^{({\rm f})}\bigr)\) &
\(\operatorname{span}_{\mathbb R}
\{\sigma_z^{({\rm N})}\otimes I_2^{({\rm f})},
\sigma_z^{({\rm N})}\otimes\sigma_x^{({\rm f})}\}\)\\
F06 &
\(\bigl(I_2^{({\rm N})}\otimes\sigma_x^{({\rm f})},
(\ii\sigma_y)^{({\rm N})}\otimes(\ii\sigma_y)^{({\rm f})}\bigr)\) &
\(\operatorname{span}_{\mathbb R}
\{\sigma_x^{({\rm N})}\otimes\sigma_y^{({\rm f})},
\sigma_z^{({\rm N})}\otimes\sigma_x^{({\rm f})}\}\)\\
F18 &
\(\bigl((\ii\sigma_y)^{({\rm N})}\otimes(\ii\sigma_y)^{({\rm f})},
I_2^{({\rm N})}\otimes\sigma_x^{({\rm f})}\bigr)\) &
\(\operatorname{span}_{\mathbb R}
\{\sigma_x^{({\rm N})}\otimes\sigma_y^{({\rm f})},
\sigma_z^{({\rm N})}\otimes\sigma_x^{({\rm f})}\}\)\\
F21 &
\(\bigl((\ii\sigma_y)^{({\rm N})}\otimes(\ii\sigma_y)^{({\rm f})},
\sigma_z^{({\rm N})}\otimes\sigma_x^{({\rm f})}\bigr)\) &
\(\operatorname{span}_{\mathbb R}
\{\sigma_y^{({\rm N})}\otimes\sigma_y^{({\rm f})},
\sigma_z^{({\rm N})}\otimes\sigma_x^{({\rm f})}\}\)
\end{tabular}
\end{ruledtabular}
\end{minipage}
\end{center}

Substitution of each row into Eq.~\eqref{eq:complete_group_space} selects
exactly the two displayed directions.  Their Fock-space lifts are the
even- and odd-parity operators in
Eqs.~\eqref{eq:odd_parity_generator} and
\eqref{eq:even_parity_generators}, up to an overall sign in either basis
direction.  The one-mode realizations and their complete space were given in
Eqs.~\eqref{eq:one_mode_realizations} and \eqref{eq:one_mode_group}.

\section{Invariant time-reversal spaces}
\label{app:antiunitary_spaces}

Expanding a general complex matrix \(U\) and solving
\(C\overline U C=U\) together with
\(\overline A\,U=UA\) for every generator gives the complete linear solution
spaces below.  The physical time-reversal matrices form the subset that also
satisfies \(U^\trans C U=C\).
\begin{equation}
\begin{array}{c|l}
\text{realization} & \text{complete linear space of invariant }U \\ \hline
\text{one-mode F01, F05, F17, F21}
&\operatorname{span}_{\mathbb R}\{I_2,\ii\sigma_z\}\\[2pt]
\text{two-mode F01}
&\operatorname{span}_{\mathbb R}
\{I_2^{({\rm N})}\otimes I_2^{({\rm f})},
I_2^{({\rm N})}\otimes\sigma_x^{({\rm f})},
\ii\sigma_z^{({\rm N})}\otimes I_2^{({\rm f})},
\ii\sigma_z^{({\rm N})}\otimes\sigma_x^{({\rm f})}\}\\[2pt]
\text{two-mode F06, F18}
&\operatorname{span}_{\mathbb R}
\{I_2^{({\rm N})}\otimes\sigma_x^{({\rm f})},
\sigma_x^{({\rm N})}\otimes\sigma_z^{({\rm f})},
\ii\sigma_y^{({\rm N})}\otimes\sigma_y^{({\rm f})},
\ii\sigma_z^{({\rm N})}\otimes I_2^{({\rm f})}\}\\[2pt]
\text{two-mode F21}
&\operatorname{span}_{\mathbb R}
\{I_2^{({\rm N})}\otimes I_2^{({\rm f})},
\sigma_x^{({\rm N})}\otimes\sigma_z^{({\rm f})},
\ii\sigma_y^{({\rm N})}\otimes\sigma_y^{({\rm f})},
\ii\sigma_z^{({\rm N})}\otimes\sigma_x^{({\rm f})}\}.
\end{array}
    \label{eq:invariant_antiunitary_spaces}
\end{equation}
Every basis matrix in Eq.~\eqref{eq:invariant_antiunitary_spaces} is symmetric.
Every real linear combination is therefore symmetric, which completes the
argument leading to Eq.~\eqref{eq:antiunitary_square}.

\bibliography{references}

\end{document}
